\subsection{The odd terminal packets with a unit arm}
\label{sec:odd-complete:q-terminals}

Let $F$ have the literal directions in
Table~\ref{tab:odd186:directions}. The two root classes considered here
are $pW_a\bar f_b$, $3\le a<b$, and $f_a h_bf_c$,
$3\le a<b<c$. Their terminal messages and their inner connector $L$
are the same. We prove the respective global fourth losses
$705/128$ and $11/2$ by comparing their complete certificates with
the auxiliary root $pP_0$ of
Lemma~\ref{lem:odd-complete:certificate}.

\paragraph{The propagated packet.}
Every terminal in this section has an actual first-row upper
\begin{equation}\label{eq:odd-q:row}
 TF\le e^\alpha(\bar f_2+BR+D),\qquad
 R=r(B)u,\quad D\ge0.
\end{equation}
For a retained square $\alpha=2\log m$; for an exponential
terminal $\alpha=-\Gamma$. Put $D_i=\E B^iR$ and
$\gamma_j=\beta_j-D_{j-2}$. The finite identity
\eqref{eq:odd-common:q-finite}, applied before discarding any term,
gives the parameters of \eqref{eq:odd-complete:input}:
\begin{equation}\label{eq:odd-q:parameters}
\begin{array}{c|c|c|c|c}
 L&C&c&Q&\E BQ(B)u\\\hline
 1&1&D&1+r&-d+D_1\\
 L\ge2&1&2D-\gamma_{L+1}&G_L+z^{L-1}r&-\gamma_{L+2}.
\end{array}
\end{equation}
The long row uses the proved nonnegative losses
$\gamma_3,\ldots,\gamma_{L+1}$, and its mean uses the additional
actual index $L+2$. Its auxiliary-root exponent is exactly
\begin{equation}\label{eq:odd-q:baseline}
 J_0=4(\alpha+c+\E BQu-d)
               +(17/4)\|(Q-B)u\|^2.
\end{equation}
In particular the constant split is $\gamma=C-1=0$ throughout.

If $r=-t+zs$ and $0\le t,s\le1$, both coordinates of every finite
$Q_L$ lie in $[0,1]$. Indeed
$Q_L=G_{L-1}+(1-t)z^{L-1}+sz^L$: each parity coordinate is an
alternating finite partial sum with its final coefficient between
zero and one. Convexity in $t,s$ reduces to the four adjacent
partial sums. This includes $Q_1=(1-t)+zs$ and
$Q_2=(1-xs)+z(1-t)$.

\paragraph{Canonical triples $(1,2,w)$.}
Use the actual square $F=qf_2h_w$. Its three retained packets are
\[
\begin{array}{c|c|c|c|c}
 w&d_{\max}&m&r&D\\\hline
 3&1/3&1-d/2&-kx&d/4\\
 4&1/4&1-d&-k(d+x)+zkx&5d/12\\
 w\ge5&1/5&1-9d/8&-ka+zkb&2d/5.
\end{array}
\qquad k=m^{-1}.
\]
The first two rows are the actual sums and signed moments in
Section~\ref{proof:Q5-C169}. The last is
Lemma~\ref{lem:reference-retained-packet} with $\sigma=9/4$;
its actual $a,b,U$ and all four residue classes of $D_i$ are
retained. In all three rows $0\le t,s\le1$, and
\eqref{eq:odd-q:baseline} has respective uniform losses
$2817/512$, $705/128$ and $6$. These are complete $v=3$
certificates: the auxiliary pair is $P_0$ in either original
direction, and its $\beta_2$ curvature coefficient is zero.

At a minimum-one root use
$\Phi(1,0,1/3)=-11255/36864<0$ and
$\Psi(1,6/5)=-8899/3840<0$.
Thus every low-domain row has loss greater than $705/128$.
On the complementary domain, the same extra arm satisfies
$h_4\le1$ for $d\ge1/4$, and $h_w\le1$ for $w\ge5$, $d\ge1/5$.
The complete message is then bounded by its actual $qf_2$
terminal. Proposition~\ref{prop:root-post:square-moments}, applied
to the one pair $(a,b)$ and terminal parameter two, supplies $705/128$.

For a positive root the $w\ge5$ low domain must extend to $1/4$.
Use instead the $\sigma=2$ row of
Lemma~\ref{lem:reference-retained-packet}: $m=1-d$,
$R=(-ka+zkb)u$, $D=d/2$, $a<7/10$, $b<3/5$.
Its own signed moments give
$\gamma_3\ge0$, $\gamma_4\ge2d/3$,
$\gamma_j\ge41d/96$ for $j\ge5$, including the mean index.
The actual short-window values of \eqref{eq:odd-q:baseline}
have losses at least $8,6,187/24,47/6$ for $L=1,2,3,4$.
For $L\ge5$ its base has modulus at most $23/16$ and its
perturbation at most $(6/5)x^2$. The latter adds less than
$2091d/800$ to the complete norm payment; the adjacent actual
losses satisfy $\gamma_{L+1}+\gamma_{L+2}\ge41d/48$.
Consequently $J_0$ has loss at least $43/6-2091/800$.
The positive-root comparison of
Lemma~\ref{lem:odd-complete:positive} saves at least $d$ on
every one of these rectangles. The weakest resulting loss is
\begin{equation}\label{eq:odd-q:retained-margin}
 49/6-2091/800=13327/2400=11/2+127/2400.
\end{equation}
For $w=3,4$ the same comparison gives $3329/512$ and $833/128$.

For the positive-root high domain use the actual reference
estimate of Proposition~\ref{prop:retained-reference-family},
$\E[f_a X^4]\le1$ for $X=T^L(qf_2h_w)$, $a\ge3$, $w\ge4$.
The remaining factor is $f_a^3(h_bf_c)^4$.
For $(a,b)=(3,4)$ the actual pair loss $69d/20$ and precursor
norm $3\sqrt d/8$ give $(h_bf_c)^4\le e^{-54d/5}$ on
$d\ge1/4$. For $b\ge5$ use $16d/5$ and $\sqrt d/4$, with
the same result. If $a=3$, then $f_3\le3/2$ and
$3\log(3/2)<5/4\le5d$; if $a\ge4$, $f_a\le1$.
The complete high-domain loss is therefore at least $29/5$.
This argument also applies to the canonical $qf_A$ reference
estimate \eqref{eq:q117-unified:shortone} for $A\ge3$ in the
following triple rows.

\paragraph{The other single triples.}
For $F=qf_Ah_C$, the actual exponential first row is
\eqref{eq:odd-q:row} with $\alpha=-\Gamma$ and $d\le1/4$.
For $A=3$,
\[
 r=-\beta_{C-1}+x^2V+zxU,
 \qquad 0\le U,V\le1,\quad |r|<3/5,\quad D=7d/20,
\]
with the true $G_1=(1,0)$ retained at $C=4$.
For $A\ge4$ use the actual four-beta packet
\[
 r=-c+(-1)^A z^A G_{C-A}(-z),\qquad
 c=\beta_{C-1}-\beta_{A-1}.
\]
Its $(\Gamma/d,D/d,\rho)$ are
$(51/16,1,3/10)$ at $A=4$,
$(89/32,2/3,1/4)$ at $A=5$ and
$(19/8,7/20,1/8)$ at $A\ge6$.
These are the same actual packets and all signed moment indices
proved in Sections~\ref{proof:Q5-C172} and~\ref{proof:Q5-C173}.
Substitution in \eqref{eq:odd-q:baseline}, with zero auxiliary
curvature, gives the following losses and positive-root increments.
An increment already includes the saved root loss.
\[
\begin{array}{c|ccc|ccc}
 &\multicolumn{3}{c|}{\text{auxiliary loss}}
 &\multicolumn{3}{c}{\text{positive-root increment upper}}\\
 A&L=1&L=2&L\ge3&L=1&L=2&L\ge3\\\hline
 3&157/20&159/20&129/20&2/3&0&0\\
 4&29/4&173/20&133/20&7/4&121/100&61/160\\
 5&57/8&1183/120&959/120&1&121/100&61/160\\
 A\ge6&41/5&231/20&191/20&5/12&121/100&61/160.
\end{array}
\]
The $A=4,L=1$ increment is strict. Here are the signed
coordinate charges that prove this table directly.
At $A=3$, $Q_1$ has real overshoot at most $x^2$;
its fee is below $119d/72$, before the saved $d$.
At $L=2$ only the imaginary coordinate may exceed one, by $x^2$;
its fee is $(17/2)(\theta^4+\theta^5)d<d$.
Lengths three and four are in the rectangle. Every later tail is
at most $x^2(1/4+x)$, below the finite partial-sum margins since
$x(5/4+x)<1$ on $[0,\theta]$.
At $A=4,L=1$, the real and negative imaginary overshoots are
at most $\delta_2+x^2$ and $x^2$, respectively. Keeping the
stronger $-11d/10$ gives the strict increment
\[
 -11/10+1547/576+17/168=34009/20160<7/4.
\]
At $A=5$ the real overshoot is only $\delta_2$, giving an
increment below $d$. At $A\ge6$ use the integral norm version
with $\rho=1/8$, giving $5d/12$.
At $L=2$ and $L\ge3$ the largest respective perturbations
$39/200$ and $39/320$ give $121d/100$ and $61d/160$.

For a minimum-one root the $A=3$ polynomials have
real coordinate nonnegative, imaginary coordinate at most $5/4$
and norm at most $3\sqrt d/2$. Their $\Phi$ cost is below $3d/5$
and their $\Psi$ cost is negative. For $A\ge4$ use instead
$\bar D=4/3$, $R=3/2$ and
$\Phi(4/3,0,1/4)=3063/4096<3/4$;
$\Psi(4/3,3/2)<0$. The low-domain losses exceed $705/128$.
In the high domain $h_C\le1$ in the unchanged direction.
At a minimum-one root apply
Proposition~\ref{prop:root-post:square-moments} with terminal
parameter $A$. At a positive root use
\eqref{eq:q117-unified:shortone} and the preceding remaining-factor
payment. The respective losses are $705/128$ and $29/5$.

\paragraph{Every odd multipair terminal of minimum one.}
Write the ordered residual list as
$(1,A_1,C_1,\ldots,A_j,C_j)$, $j\ge2$.
The canonical direction is $qf_{A_1}h_{C_1}\prod_{i\ge2}h_{A_i}f_{C_i}$.
When $(A_1,C_1,A_2)=(2,3,4)$ use the fixed cross
$qf_2h_{C_2}h_3f_4\prod_{i\ge3}h_{A_i}f_{C_i}$.
The single finite mask and its actual first-row moments are those
of Section~\ref{proof:Q5-C177}. The five records below give
$d_*$, $\Gamma/d$, $D/d$, a modulus bound $\rho$ for $r$,
and the complete auxiliary losses at $L=1,2,\ge3$.
\[
\begin{array}{c|c|c|c|c|ccc}
 &d_*&\Gamma/d&D/d&\rho&1&2&\ge3\\\hline
 \mathrm I:A_1\ge3&1/64&431/80&5/3&1/2&1333/160&659/96&3799/480\\
 \mathrm {II}:A_1=2,C_1\ge4&1/64&351/80&11/10&2/3&6287/480&1027/96&715/96\\
 \mathrm {III}:A_1=2,C_1=3,A_2\ge5&1/8&21/5&11/10&3/5&719/60&115/12&47/6\\
 \mathrm {IV}:\text{cross},j=2&1/8&21/4&11/10&1/2&971/60&827/60&757/60\\
 \mathrm V:\text{cross},j\ge3&1/8&169/20&7/4&3/5&953/60&983/60&589/30.
\end{array}
\]
The table conservatively keeps the original uniform root-curvature
fees $5d/96$ or $5d/12$; their actual auxiliary values are zero.
The same four residue classes prove all finite $\gamma$ signs.

For the minimum-one transfer every $Q_L$ has nonnegative real
coordinate, imaginary coordinate at most $3/2$ and norm at most
$8\sqrt d/5$. Indeed the actual $r=A_r+zD_r$ satisfy
$-2/3\le A_r\le1/2$, $|D_r|\le2/3$.
For $L=1,2$ use the explicit rows $(1+A_r,D_r)$ and
$(1-xD_r,1+A_r)$; for $L\ge3$ use
$1-x-\rho x>0$, $1+\rho\sqrt x<3/2$ and
$|Q_L|\le1+\rho x<8/5$. Thus $\Phi$ costs less than
$11d/10$ in I--II and $6d/5$ in III--V, and $\Psi$ costs
less than $7d/10$ throughout. Every residual loss exceeds
$705/128$, for example $659/96-11/10=2767/480$.

The positive-root increments are controlled by the same signed
overshoot functional, as follows:
\[
\begin{array}{c|ccc}
 &L=1&L=2&L\ge3\\\hline
 \mathrm I&1547/600+17/112-1&4/5&125/96\\
 \mathrm {II}&-1&-1&-1\\
 \mathrm {III}&0&0&141/80\\
 \mathrm {IV}&0&0&125/96\\
 \mathrm V&29/5&171/50&141/80.
\end{array}
\]
For I, $A_r\le13/50$ and $D_r\in[-1/4,2/3]$.
At $L=1$ these give the two fees $1547d/600$ and $17d/112$;
at $L=2$ the real and imaginary overshoots are at most $x/4$
and $13/50$, giving a fee below $9d/5$ before saving $d$.
For II, $r=-t+zs$ lies in $[0,2/3]^2$ in its signed coordinates,
so every finite row is in the rectangle.
For III, $0\le t\le2/3$, $|s|\le1/5$; the first two fees
are $17d/140$ and $119d/144$, both below the saved $d$.
For IV, $0\le t\le19/40$, $-x^2\le s\le19/128$;
the corresponding fees are $17d/168$ and $17d/24$.
For longer rows of I, III and IV use the integral norm estimate
with $\rho x$. For V use that estimate at every length with
$\rho=3/5$, $\sqrt\theta<13/20$ and $\theta=13/32$.
Each entry subtracted from its auxiliary loss leaves more than six.

For $d\ge1/64$ every later actual pair with minimum at least
four is at most $e^{-88d/125}\le1$. Hence I--II compare with
their first triple. On $d\ge1/8$, III can be written as the
same product
\[
 F=(qf_2h_{A_2})(h_3f_{C_2})\prod_{i\ge3}h_{A_i}f_{C_i};
\]
IV--V already have the analogous factors $qf_2h_{C_2}$ and
$h_3f_4$. The actual centered-row proof gives, on this domain,
\begin{equation}\label{eq:odd-q:high-three-pair}
 h_3f_t\le e^{-d},\qquad f_3h_t\le e^{-d},\qquad t\ge4.
\end{equation}
For $t\ge5$ its centered fluctuation has absolute value at most
$\sqrt{a_0^2+\operatorname{Var}U}\le a_0+d$;
here $U=\sum_{r=3}^{t-1}K^r\one$ and
$a_0=\beta_{t-1}-\E U=\Gamma_t-2d\ge0$, with
$\Gamma_t=d+\beta_t+\beta_{t-1}$.
The centered original row gives this bound for
$a_0g-K(U-\E U)$, $g=-u$, since its two row coordinates satisfy
$g^2+pq=1$. The finite $t=5$ row uses
$8\theta^2(1+\theta)<35/16$, and the $t\ge6$ tail uses
$\operatorname{Var}U\le\theta^2d\le11d^2/8$.
At four use $\sqrt{d^2+\delta_3}\le2d$ in the exact arm sum.
The two literal sums have opposite signs of this same centered
fluctuation, so its absolute-value bound proves both inequalities.
The resulting actual comparison is with the canonical triple
already proved, in the same arm directions. This gives the
claimed global losses in both root classes.

All parameter intervals include their nonzero endpoints.
At $d=0$, every vector here vanishes and leaf integration gives
the value one.
